Fix a class member \(P\), abbreviate \(\mu=\mu_P\) and \(\pi=\pi_P\), and write
\[
f(a)=\int\pi(a\mid x)\,dP_X(x),
\qquad
\theta(a)=\int\mu(a,x)\,dP_X(x).
\]
Because \(P_X\) is a probability law and \(c\le\pi\le c^{-1}\),
\[
c\le f(a)\le c^{-1}\quad(0\le a\le1).
\]
Integration of the uniform treatment-direction Hölder inequalities against \(P_X\), justified by boundedness, shows that \(\theta\) is univariate Hölder of order \(\alpha\). Thus, with \(p=\lfloor\alpha\rfloor\), there is for every sufficiently small \(h\) a polynomial \(Q_h\) of degree at most \(p\), centered at \(t_0\), such that
\[
Q_h(t_0)=\theta(t_0),
\qquad
\sup_{|a-t_0|\le h}|\theta(a)-Q_h(a)|\le Ch^\alpha.
\]
At integer \(\alpha\), the Taylor polynomial through order \(\alpha-1\) gives this bound and is contained in the degree-\(p\) space.

Let \(S(u)=(1,u,\ldots,u^p)^\top\), \(S_h(a)=S((a-t_0)/h)\), and \(K_h(a)=h^{-1}K((a-t_0)/h)\). The population marginal Gram is
\[
D_h=\int K(u)S(u)S(u)^\top f(t_0+hu)\,du
\succeq c\int K(u)S(u)S(u)^\top\,du.
\]
The last fixed matrix is positive definite, since \(K\) is bounded below on an interval and a nonzero polynomial cannot vanish on that interval. Hence \(\lambda_{\min}(D_h)\ge\lambda_0>0\). For every polynomial \(q\) of degree at most \(p\), the equivalent population weight
\[
W_h(a)=e_0^\top D_h^{-1}S_h(a)K_h(a)
\]
satisfies the exact reproduction identity
\[
\int W_h(a)q(a)f(a)\,da=q(t_0).
\]
Moreover, \(\int|W_h(a)|f(a)\,da\le C\). Therefore the population local-polynomial intercept of any response with conditional mean \(\theta(a)\) has bias at most \(Ch^\alpha\), without any derivative or Hölder modulus of \(f\).

We next establish the nuisance remainder. Let \(\widetilde\mu(a,x)=\widehat\mu_0(x)\), \(\widetilde\pi(a\mid x)=\widehat\pi_0(x)\), and use population marginalizations temporarily. Put
\[
v_a(x)=\widetilde\mu(a,x)-\mu(a,x),
\qquad
e_a(x)=\widetilde\pi(a\mid x)-\pi(a\mid x),
\qquad
\overline e_a=P_Xe_a.
\]
Bayes' formula gives the conditional law of \(X\) given \(A=a\) density \(\pi(a\mid x)/f(a)\) with respect to \(P_X\). Direct conditioning of
\[
\widetilde\xi=
\frac{P_X\widetilde\pi(A\mid\cdot)}{\widetilde\pi(A\mid X)}
\{Y-\widetilde\mu(A,X)\}+P_X\widetilde\mu(A,\cdot)
\]
therefore yields the exact identity
\[
\mathbb E_P[\widetilde\xi\mid A=a]-\theta(a)
=P_X\left[\frac{v_ae_a}{\widetilde\pi_a}\right]
-\frac{\overline e_a}{f(a)}P_X\left[\frac{v_a\pi_a}{\widetilde\pi_a}\right]
=:B(a).
\]
Since \(c/2\le\widetilde\pi\le2/c\), Cauchy--Schwarz gives
\[
|B(a)|\le C\lVert v_a\rVert_{L_2(P_X)}\lVert e_a\rVert_{L_2(P_X)}.
\]
Let
\[
U=\lVert\widehat\mu_0-\mu(t_0,\cdot)\rVert_2,
\qquad
V=\lVert\widehat\pi_0-\pi(t_0\mid\cdot)\rVert_2.
\]
The treatment Hölder conditions and the constant-in-dose extensions imply, for \(|a-t_0|\le h\),
\[
\lVert v_a\rVert_2\le U+Lh^{A_0},
\qquad
\lVert e_a\rVert_2\le V+Lh^{B_0}.
\]
Consequently,
\[
\sup_{|a-t_0|\le h}|B(a)|
\le C\{UV+Uh^{B_0}+Vh^{A_0}+h^{A_0+B_0}\}
=C(U+h^{A_0})(V+h^{B_0}).
\]
The term \(h^{A_0+B_0}\) is retained: it is a residual of this constant-in-dose first-order construction, not a claimed lower bound for the model class.

By \(\mathrm{lem:separate\text{-}slice\text{-}pilot\text{-}mse}\), \(\mathbb E U^2\le Cu^2\) and \(\mathbb E V^2\le Cv^2\). Crucially, \(\mathrm{lem:independent\text{-}pilot\text{-}product}\) gives
\[
\mathbb E[U^2V^2]=\mathbb E[U^2]\mathbb E[V^2]\le Cu^2v^2.
\]
Expanding the square of the preceding remainder bound, using boundedness for the purely linear cross terms, and applying these three second-moment bounds gives
\[
\mathbb E\left[\sup_{|a-t_0|\le h}|B(a)|^2\right]
\le C\{(u+h^{A_0})(v+h^{B_0})\}^2.
\]
No fourth moment of an integrated pilot norm is invoked.

It remains to pass to the empirical marginalizations and empirical local-polynomial fit. Conditional on both pilot folds, the two empirical marginalization errors on \(I_X\) have conditional second moments bounded by \(C/n\), because all clipped pilot functions are bounded. Substitution in the displayed exact remainder identity shows that they add at most \(C/n\) to mean-squared error; all multiplying factors are bounded by positivity and clipping.

For the empirical Gram \(\widehat D_h\), each centered entry is an average of independent variables bounded by \(C/h\) with second moment at most \(C/h\). The elementary exponential inequality
\[
e^x\le1+x+\frac{x^2}{2(1-|x|/3)},\qquad |x|<3,
\]
followed by exponential Markov optimization gives, entrywise, a Bernstein bound of order \(\exp(-cnh)\) for a fixed deviation. The Gram dimension \(p+1\) is fixed, so a union bound and the fixed-dimensional operator-norm inequality yield
\[
P\{\lambda_{\min}(\widehat D_h)<\lambda_*\}\le Ce^{-c_1nh}.
\]
On the complementary event, \(\lVert\widehat D_h^{-1}\rVert\le\lambda_*^{-1}\), and empirical polynomial reproduction is exact. Applying it to \(Q_h\), the Taylor remainder has root mean square at most \(Ch^\alpha\), since
\[
\mathbb E\left[\left\{N_{\mathrm{lp}}^{-1}\sum_{i\in I_{\mathrm{lp}}}|K_h(A_i)|\right\}^2\right]
\le C\left(1+\frac1{nh}\right).
\]
Conditional on the training folds and treatment values, the centered pseudo-outcome score has variance bounded by
\[
\frac{C}{N_{\mathrm{lp}}^2}\sum_{i\in I_{\mathrm{lp}}}K_h(A_i)^2\lVert S_h(A_i)\rVert^2,
\]
whose expectation is at most \(C/(nh)\). Combining polynomial bias, nuisance remainder, empirical marginalization, and conditional variance proves on the good-Gram event
\[
\mathbb E_P[(\widehat\theta-\theta(t_0))^2;\,\mathrm{good}]
\le C\left[h^{2\alpha}+\{(u+h^{A_0})(v+h^{B_0})\}^2+\frac1{nh}+\frac1n\right].
\]
The final clipping is nonexpansive relative to \(\theta(t_0)\in[-M_Y,M_Y]\). On the bad-Gram event, both the clipped estimate and target lie in that interval, so the risk contribution is at most \(4M_Y^2Ce^{-c_1nh}\). All operations in the construction are Borel measurable. The constants used above are uniform over the declared class, and taking the supremum over \(P\) proves the result.